% This file was converted to LaTeX by Writer2LaTeX ver. 1.9.9
% see http://writer2latex.sourceforge.net for more info
\documentclass{article}
\usepackage{calc,amsmath,amssymb,amsfonts}
\usepackage[T1]{fontenc}
\usepackage[italian]{babel}
\usepackage[style=numeric,backend=biber]{biblatex}
\author{HP}
\date{2026-07-26}
\begin{document}
\section{Prefazione}
Ci sono suoni che non si dimenticano.

Il fruscio del televisore mentre cercava il canale giusto. Il clic deciso dell'interruttore di alimentazione. Il rumore
inconfondibile del Datasette che iniziava a girare. E poi quell'attesa. Lunga. A volte interminabile. Un'attesa fatta
di speranza, perché bastava un piccolo errore, un nastro leggermente rovinato o un colpo involontario al registratore
per costringerti a ricominciare tutto da capo.

Se stai leggendo queste pagine, probabilmente sai esattamente di cosa stiamo parlando.

Forse eri uno di quei ragazzi che, negli anni Ottanta, passavano interi pomeriggi davanti a un televisore trasformato in
monitor. Magari avevi il naso incollato al manuale del Commodore 64, cercando di capire cosa significasse davvero una
riga di BASIC. Oppure aspettavi con impazienza l'uscita della tua rivista preferita, pronto a digitare centinaia di
righe di codice senza sapere se, alla fine, avresti visto comparire un'astronave, una pallina rimbalzante o
semplicemente un messaggio di errore.

Eppure non importava.

Perché ogni tentativo era una scoperta.

Ogni errore insegnava qualcosa.

Ogni programma funzionante sembrava una piccola magia.

Il Commodore 64 non è stato soltanto uno dei computer più importanti della storia dell'informatica. Per molti di noi è
stato il primo insegnante. Ci ha mostrato che un computer non era una macchina misteriosa costruita da qualcun altro.
Era qualcosa con cui si poteva dialogare. Bastava impararne il linguaggio.

Quarant'anni dopo, molte cose sono cambiate.

I computer sono diventati milioni di volte più potenti. Le connessioni sono istantanee. I linguaggi di programmazione
sono evoluti. Gli strumenti di sviluppo fanno in pochi secondi ciò che allora richiedeva giornate intere.

E poi è arrivata l'intelligenza artificiale.

Quando abbiamo iniziato a immaginare questo libro ci siamo posti una domanda molto semplice.

Che cosa succederebbe se quel ragazzo degli anni Ottanta tornasse oggi davanti al suo Commodore 64?

La risposta non era quella che ci aspettavamo.

Non sarebbe cambiato il computer.

Sarebbe cambiato il compagno di laboratorio.

Negli anni Ottanta c'era sempre qualcuno che ne sapeva un po' più di noi. L'amico del quartiere che aveva già imparato
un comando nuovo. Il compagno di scuola che aveva scoperto un POKE misterioso. Il negoziante che consigliava una
rivista. L'autore di un listato pubblicato il mese precedente.

Oggi quel compagno di laboratorio può assumere una forma diversa.

Può chiamarsi ChatGPT.

Questo libro nasce proprio da quell'incontro.

Non è stato scritto da un essere umano che ha chiesto a un'intelligenza artificiale di fare il lavoro al suo posto.

E non è nemmeno un testo prodotto da un'intelligenza artificiale.

È il risultato di un dialogo.

Di centinaia di domande.

Di confronti.

Di idee migliorate insieme.

Di capitoli riscritti più volte.

Di esperimenti.

Di entusiasmi condivisi.

Abbiamo deciso di rendere questo processo parte integrante del libro perché crediamo che rappresenti un nuovo modo di
imparare. L'intelligenza artificiale non sostituisce la curiosità. Non elimina gli errori. Non cancella la
soddisfazione della scoperta. Può suggerire una strada, ma sarà sempre il lettore a percorrerla.

Per questo motivo il vero protagonista di queste pagine non è il Commodore 64.

Non è nemmeno ChatGPT.

Il protagonista sei tu.

Tu che, capitolo dopo capitolo, costruirai un videogioco funzionante.

Lo faremo insieme, partendo dalle basi del linguaggio BASIC, entrando poi nel mondo dell'Assembly, imparando a conoscere
il VIC-II, il SID, gli sprite, gli interrupt raster, lo scrolling, la memoria e tutte quelle tecniche che hanno
permesso a una macchina con appena 64 kilobyte di RAM di entrare nella storia.

Ma questo non sarà un manuale nel senso tradizionale del termine.

Non troverai una lunga sequenza di definizioni da memorizzare.

Ogni argomento nascerà perché il nostro videogioco ne avrà bisogno.

Quando il protagonista dovrà comparire sullo schermo, scopriremo gli sprite.

Quando il gioco inizierà a rallentare, capiremo perché il BASIC non basta più.

Quando vorremo rendere tutto più fluido, entreremo nel mondo dell'Assembly.

Ogni capitolo aggiungerà davvero qualcosa al progetto.

Ogni nuovo concetto renderà il videogioco un po' più completo.

Alla fine del percorso non avrai soltanto letto un libro.

Avrai costruito qualcosa che prima non esisteva.

Forse, mentre programmeremo insieme, ti tornerà in mente qualche ricordo. Il pomeriggio trascorso a copiare un listato.
La gioia del primo programma funzionante. L'emozione di vedere uno sprite attraversare lo schermo per la prima volta.
Oppure scoprirai tutto questo oggi, magari perché il Commodore 64 appartiene a un'epoca che non hai vissuto, ma che
desideri conoscere.

In entrambi i casi, il nostro obiettivo sarà lo stesso.

Non insegnarti semplicemente a programmare un Commodore 64.

Vorremmo accompagnarti in un viaggio.

Un viaggio verso un tempo in cui ogni riga di codice era una scoperta, ogni byte aveva un valore e ogni piccolo successo
sembrava aprire una porta su un mondo nuovo.

Se, arrivato all'ultima pagina, guarderai ancora una volta quella schermata blu con il cursore lampeggiante e sentirai
nascere il desiderio di scrivere un programma tutto tuo, allora questo libro avrà raggiunto il suo scopo.

Perché il vero traguardo non sarà il videogioco che avremo costruito insieme.

Sarà aver riscoperto il piacere di essere curiosi.


\bigskip
\end{document}
